\begin{textsample}{NEG Dim 2 – vem\_ed\_34 – Score 3.00 – vem\_ed\_34/t184.txt}  \label{ex:f2_neg_064}
Tecendo conexões Realizado no segundo semestre de 2025, o Projeto Colaborativo Internacional “Conociendo las características del tejido denim en Perú y Brasil: Diferencias y similitudes en los procesos de desarrollo y lavados”, realizado entre a Fatec Americana e o Centro de Altos Estudios de la Moda (CEAM, Peru), teve como objetivo principal proporcionar aos participantes uma imersão comparativa nas particularidades do tecido denim, abrangendo desde suas características técnicas e estéticas até os processos de desenvolvimento e beneficiamento, com foco nas realidades produtivas do Brasil e do Peru.
A iniciativa visou fomentar a colaboração internacional, o intercâmbio cultural e o desenvolvimento de competências essenciais para o mercado global da moda.
Tanto a Fatec Americana quanto o CEAM têm cursos focados na área de moda e têxtil, o que permitiu um intercâmbio de conhecimentos sobre processos e tecnologias específicas do setor. “Foi uma excelente oportunidade para os alunos conhecerem uma nova cultura, de maneira geral; e dentro da própria disciplina, entender como os processos de beneficiamento do denim são feitos em outro país.
É um conhecimento que vai além da sala de aula”, comenta Daives Araken Bergamasco, que conduziu o PCI junto ao curso de Design de Moda, em colaboração com a professora peruana Aida Esperanza Custodio Vivar de Valencia.
Daives ressalta a necessidade de desenvolver uma boa estrutura de projeto, mas com flexibilidade para fazer as adaptações que forem necessárias.
Destaca ainda a importância de estimular os estudantes para que a interação com os colegas estrangeiros não seja “robótica”.
O professor relata que “no início, a comunicação foi bem tímida, mas do meio para o final do projeto, os alunos conseguiram estabelecer uma relação muito sólida e focada nas atribuições dos trabalhos propostos”.
Os estudantes interagiram em grupos de WhatsApp e publicaram os resultados dos trabalhos no Padlet.
Os peruanos elaboraram os croquis das calças jeans e os brasileiros escolheram os melhores tecidos e processos produtivos para executar as peças.
O aluno Leonardo Lopes de Medeiros resume entusiasmado sua participação no projeto: “Participar do PCI / COIL com a faculdade CEAM, do Peru, foi uma das experiências mais marcantes da minha graduação.
Trocar ideias com estudantes de outro país, com outra cultura e outra forma de ver o mundo da moda, abriu minha cabeça de um jeito que nenhuma aula convencional conseguiria.
Foi incrível perceber como os processos criativos podem ser tão diferentes dependendo de onde você vive.
A gente conseguiu comparar nossas referências, aprender com as diferenças e, juntos, mostrar como transformar ideias criativas em produtos reais de moda — o que, no fim das contas, é exatamente o que a gente estuda para fazer.
Além de tudo que aprendi na prática, essa experiência me mostrou o quanto o setor têxtil é relevante globalmente e como a colaboração entre países pode fazer a formação da gente ser muito mais rica.
O projeto prova que aprender vai além das paredes da sala de aula — e que, quando a gente abre espaço para o diferente, o resultado é sempre mais interessante.”

% matched lemmas: 
\end{textsample}
